% Generated by src/equivalent_rights_proposals.jl -- do not edit by hand.
\begin{table}[htbp]
\centering
\small
\caption{Club-price rights ratios $\rho_i$ equivalent to the Wolfram et al. and Banerjee, Duflo \& Greenstone proposals}
\renewcommand{\arraystretch}{1.15}
\resizebox{\textwidth}{!}{
\begin{tabular}{lccccc}
  \toprule
  & \multicolumn{2}{c}{\textbf{Wolfram et al.}} & \multicolumn{3}{c}{\textbf{Banerjee, Duflo \& Greenstone}} \\
  \cmidrule(lr){2-3} \cmidrule(lr){4-6}
  \textbf{Country} & $p_i$ & $\rho_i$ & $\hat\rho_i$ & $p_i$ & $\rho_i$ \\
  \midrule
  United States & 0 & -- & 3.08 & 75 & 2.83 \\
  Russia & 0 & -- & 1.99 & 75 & 1.79 \\
  European Union & 75 & 0.47 & 0.42 & 75 & 0.45 \\
  Turkey & 0 & -- & 1.5 & 50 & 1.13 \\
  China & 50 & 0.81 & 1.04 & 50 & 0.82 \\
  Indonesia & 50 & 0.72 & 1.05 & 50 & 0.77 \\
  Brazil & 50 & 0.27 & 0.48 & 50 & 0.29 \\
  India & 25 & 0.59 & 1.16 & 30 & 0.62 \\
  Nigeria & 0 & -- & 0.24 & 30 & 0.12 \\
  DR Congo & 0 & -- & 0.09 & 10 & 0.03 \\
  \midrule
  \multicolumn{6}{l}{\textit{Reduced emissions (variant 1)}} \\
  \textbf{\quad Emissions change in the coalition (\%)} &  & $-$34.5\% &  &  & $-$27.2\% \\
  \textbf{\quad World temp.~2100, change (\textdegree{}C)} &  & $-$0.068 &  &  & $-$0.108 \\
  \textbf{\quad World welfare gain (\%)} &  & 0.322\% &  &  & 0.618\% \\
  \textbf{\quad World consumption gain (\%)} &  & $-$0.012\% &  &  & $-$0.020\% \\
  \multicolumn{6}{l}{\textit{Increased consumption (variant 2)}} \\
  \textbf{\quad World welfare gain (\%)} &  & 0.025\% &  &  & 0.224\% \\
  \textbf{\quad World consumption gain (\%)} &  & 0.016\% &  &  & 0.045\% \\
  \multicolumn{6}{l}{\textit{Surplus shared by marginal utility (variant 3)}} \\
  \textbf{\quad World welfare gain (\%)} &  & 0.650\% &  &  & 1.323\% \\
  \textbf{\quad World consumption gain (\%)} &  & 0.031\% &  &  & 0.055\% \\
  \multicolumn{6}{l}{\textit{Surplus shared by uniform scaling (variant 4)}} \\
  \textbf{\quad World welfare gain (\%)} &  & $-$0.057\% &  &  & 0.063\% \\
  \textbf{\quad World consumption gain (\%)} &  & 0.013\% &  &  & 0.040\% \\
  \bottomrule
\end{tabular}
}
\label{tab:equiv_rights_ede}
\\[4pt]
{\footnotesize Note: $p_i$ is the carbon price the proposal asks of country $i$ in 2030; the comparison regime instead prices every member of the proposal's own club at one common price (\$49.2/t for Wolfram et al. and \$55.8/t for Banerjee, Duflo \& Greenstone in 2030) and leaves non-members unpriced, as the proposal does. $\rho_i$ is the country's allocation as a multiple of an equal-per-capita share of the club's emissions under the proposal, taken at face value in variant 1; $\rho^{\mathrm{isolated}}$ solves each country on its own, $\rho^{\mathrm{joint}}$ all of them jointly. $\hat\rho_i$ is the predicted $\rho_i$ for the Banerjee, Duflo \& Greenstone schedule: the country's own emissions under the proposal over an equal-per-capita share of the club's, each year weighted by the discounted club price over 2025--2100. The European Union figure aggregates its members' rights, and a country the proposal does not price is outside the club in both regimes, shown at $p_i=0$ with no equivalent allocation (`--'). Variant 1 (reduced emissions) lets the equivalent allocation set the cap: in the joint solve every member is held exactly at its level under the proposal; in the isolated solve each country's ratio is solved with all others grandfathered, and the ratios are then applied together. Variants 2 to 4 (increased consumption) keep club emissions at the proposal's and share the surplus: equalising the relative gains of all members (joint variant 2), by scaling the isolated allocation up uniformly (isolated variant 2, also variant 4), and in proportion to population times marginal utility (variant 3). The temperature row is the change from the proposal's own 2100 warming, which is 2.16\textdegree{}C for Wolfram et al. and 1.80\textdegree{}C for Banerjee, Duflo \& Greenstone. World welfare is the NPV of the inequality-averse global EDE consumption ($\eta=1.5$); world consumption is mean consumption per capita, which weights every person's consumption equally.}
\end{table}
